% R589: EGA I, Chapter0, whole3.2.3-5; source-faithful translation.
% Traducción: OpenAI Codex, GPT-6.1 Sol, esfuerzo Ultra.
\begin{env}[3.2.3]
\label{0.3.2.3-fr}
Sean $\mathcal{F}$ y $\mathcal{G}$ dos prehaces sobre la base
$\mathfrak{B}$, con valores en $\boldsymbol{K}$; se define un \emph{morfismo}
$u:\mathcal{F}\to\mathcal{G}$ como una familia
$(u_V)_{V\in\mathfrak{B}}$ de morfismos
$u_V:\mathcal{F}(V)\to\mathcal{G}(V)$ que satisfacen las condiciones
usuales de compatibilidad con los morfismos de restricción $\rho_V^W$. Con
la notación de (3.2.1), se deduce un morfismo
$u':\mathcal{F}'\to\mathcal{G}'$ de los prehaces (ordinarios) correspondientes
tomando como $u_U'$ el límite proyectivo de los $u_V$ para
$V\in\mathfrak{B}$ y $V\subset U$; la verificación de las condiciones de
compatibilidad con los $\rho'_U{}^{U'}$ se deduce del carácter funtorial del
límite proyectivo.
\end{env}

\begin{env}[3.2.4]
\label{0.3.2.4-fr}
Si la categoría $\boldsymbol{K}$ admite límites inductivos y $\mathcal{F}$
es un prehaz sobre la base $\mathfrak{B}$, con valores en $\boldsymbol{K}$,
para todo $x\in X$ las vecindades de $x$ que pertenecen a $\mathfrak{B}$ forman
un conjunto cofinal (para $\supset$) en el conjunto de las vecindades de $x$;
por tanto, si $\mathcal{F}'$ es el prehaz (ordinario) correspondiente a
$\mathcal{F}$, el tallo $\mathcal{F}_x'$ es igual a
\[
  \varinjlim_{\mathfrak{B}}\mathcal{F}(V)
\]
tomado sobre el conjunto de los $V\in\mathfrak{B}$ que contienen a $x$.
Si $u:\mathcal{F}\to\mathcal{G}$ es un morfismo de prehaces sobre
$\mathfrak{B}$ con valores en $\boldsymbol{K}$ y
$u':\mathcal{F}'\to\mathcal{G}'$ es el morfismo correspondiente de prehaces
ordinarios, $u_x'$ es, del mismo modo, el límite inductivo de los morfismos
$u_V:\mathcal{F}(V)\to\mathcal{G}(V)$ para $V\in\mathfrak{B}$ y $x\in V$.
\end{env}

\begin{env}[3.2.5]
\label{0.3.2.5-fr}
Volvamos a las condiciones generales de (3.2.1). Si $\mathcal{F}$ es un
\emph{haz} ordinario con valores en $\boldsymbol{K}$ y $\mathcal{F}_1$ es el
haz \emph{sobre $\mathfrak{B}$} obtenido al restringir $\mathcal{F}$ a
$\mathfrak{B}$, el haz ordinario $\mathcal{F}_1'$ obtenido a partir de
$\mathcal{F}_1$ mediante el procedimiento de (3.2.1) es canónicamente isomorfo a
$\mathcal{F}$, por la condición (F) y las propiedades de unicidad del
límite proyectivo. Habitualmente se identificarán $\mathcal{F}$ y
$\mathcal{F}_1'$.

Si $\mathcal{G}$ es un segundo haz (ordinario) sobre $X$ con valores en
$\boldsymbol{K}$ y $u:\mathcal{F}\to\mathcal{G}$ es un morfismo, la observación
anterior muestra que los datos de los
$u_V:\mathcal{F}(V)\to\mathcal{G}(V)$ \emph{solo para los $V\in\mathfrak{B}$}
determinan completamente $u$; recíprocamente, dados los $u_V$
para $V\in\mathfrak{B}$, basta verificar la conmutatividad del diagrama
con los morfismos de restricción $\rho_V^W$ para $V\in\mathfrak{B}$,
$W\in\mathfrak{B}$ y $V\subset W$, para que exista un morfismo $u'$ y uno
solo de $\mathcal{F}$ a $\mathcal{G}$ tal que $u_V'=u_V$ para todo
$V\in\mathfrak{B}$ (3.2.3).
\end{env}
